\subsection{Evidence and measurement design}\label{supp:prediction-evidence}
The detailed matrix supports the compact experimental tests in article
Table~\ref{tab:predictions}. Model comparisons and biological observations
are distinguished; the latter do not calibrate the detector.
\begingroup\small
\renewcommand{\arraystretch}{1.15}
\begin{longtable}{@{}>{\raggedright\arraybackslash}p{0.30\textwidth}
 >{\raggedright\arraybackslash}p{0.27\textwidth}
 >{\raggedright\arraybackslash}p{0.35\textwidth}@{}}
\caption{Predictions, existing evidence, and tests. Published observations provide
qualitative comparisons; their quantitative reanalysis remains to be done.}
\label{tab:prediction-evidence}\\
\toprule
\textbf{Prediction (model basis)} & \textbf{Existing evidence} & \textbf{Proposed test} \\
\midrule\endfirsthead
\multicolumn{3}{@{}l}{\textit{Table \thetable\ continued}}\\
\toprule
\textbf{Prediction (model basis)} & \textbf{Existing evidence} & \textbf{Proposed test} \\
\midrule\endhead
\midrule\multicolumn{3}{r@{}}{\textit{Continued on next page}}\\\endfoot
\bottomrule\endlastfoot
\textbf{Carrier dependence.} At fixed local polarization, the leading beat is
carrier-independent in the static limit and fast-carrier QIF asymptote, provided
no other circuit mode is excited. Access restores dependence at fixed applied
field; detector dynamics can add it near the crossover.
& Network sensitivity, cellular mixing, and thresholds depend on carrier frequency
\cite{esmaeilpour2021,luff2024mixer,opancar2025,peripheral2023}; bare kHz fields
have limited effects on hippocampal oscillations \cite{esmaeilpour2020}.
Local polarizations were not held fixed.
& Measure each local polarization while sweeping $f_c$ at fixed beat and state;
normalize the weak beat by their product. Hold network gain fixed or measure it
independently to separate access from residual detector dependence. Measure any
residual cutoff rather than assigning it to a compartment. \\[4pt]
\textbf{Independent carrier amplitudes.} At fixed detector and network gains,
quadratic modulation excursion follows $2|\varepsilon_1\varepsilon_2|$;
an ideal envelope follower tracks $2\min(|\varepsilon_1|,|\varepsilon_2|)$
(Eq.~\eqref{eq:reweight}, Table~\ref{tab:transduction-maps}).
& The unequal-amplitude spatial example separates these maps; a biological test
with independent local-polarization calibration is not supplied here.
& Hold the weaker polarization fixed and vary the stronger, then swap tones.
Keep both inputs weak, frequencies fixed, and the operating state controlled.
Reconstruct the slow-response excursion after correcting network filtering;
a beat lock-in alone does not measure the envelope's full excursion. \\[4pt]
\textbf{Current-ratio steering.} With fixed electrode pairs, frequencies,
access, and neural state, the quadratic beat map scales with
$I_1^{\rm pair}I_2^{\rm pair}$; its normalized pattern does not move.
An envelope follower can instead shift its maximum.
& The product-law restriction \cite{cassara2025part2} holds in IXI002
(\S\ref{sec:fem-results}); a physiological beat-map test remains open.
& Vary the current ratio at fixed electrodes and compare normalized beat maps
and target-to-reference ratios, using a cortex-only reference to test deep-target
selectivity. Keep the projection direction fixed at each site,
including any direction chosen to maximize the envelope at a reference ratio.
Measure the DC shift and repeat at lower amplitude to check state dependence.
 \\[4pt]
\textbf{State and carrier dependence.} The static sigmoid follows $\Sig''$ and
reverses sign at its inflection (PEIX$=0$). Fast-carrier QIF mixes in voltage and
reads the source through rate and network susceptibility; its DC term follows
$F'$ without the same sign rule (\S\ref{sec:dynamic}).
& Cell-type-specific responses (PV versus pyramidal) are reported
\cite{caldas2024}; the detector coefficient remains model-dependent.
& Vary tonic drive and carrier frequency at fixed beat; measure beat amplitude
and phase. Compare with the predicted coefficient after accounting for network
susceptibility. Amplitude changes alone do not isolate detection. \\[4pt]
\textbf{Transduction class.} At a matched site, weak smooth detection with nonzero
curvature gives $Y_\Omega\propto\varepsilon^2$ and a harmonic ratio vanishing as
$\varepsilon\to0$. An ideal envelope follower gives $Y_\Omega\propto\varepsilon$
and ratio $1/5$; a threshold has onset $2\varepsilon=\delta$ and a variable ratio.
& Untested with local-polarization calibration; existing comparisons match
extracellular field rather than $\varepsilon$.
& Sweep amplitude off resonance, reading both the log--log slope and
$|Y_{2\Omega}/Y_\Omega|$. Near resonance, suppression by approximately
$3\mathcal Q$ requires correction with an independently measured
$G(2\Omega)/G(\Omega)$ (\S\ref{sec:transduction-methods}). \\[4pt]
\textbf{Frequency selectivity.} The beat response peaks at the region's resonant
rhythm and grows as $1/\gamma$ near its Hopf; state changes can alter both the
generated drive and the recurrent gain \cite{ruffini2025cfc}.
& The quantitative gain--bandwidth relation is not tested here in biological data.
& Sweep $\Delta f$ through the intrinsic band across brain states and track the
peak frequency, height, and width. \\[4pt]
\textbf{Population modulation and mean rate.} Beat modulation can exceed the mean shift; their ratio
peaks and then falls as the amplified oscillation rectifies
(\S\ref{sec:timing}--\ref{sec:qif}).
& Qualitatively consistent with TI altering spike timing more than mean rate in
primates \cite{vieira2024}.
& Compare beat lock-in with mean-rate change in commensurate effect measures,
retaining the finite-amplitude qualification. \\[4pt]
\textbf{Network readout.} Population nonlinearity and synaptic recurrence contribute
to the measured response.
& Most pyramidal TI responses vanish under synaptic blockade \cite{caldas2024},
supporting network involvement without excluding cellular mixing.
& Test whether synaptic blockade reduces recurrent amplification while preserving
a cellular detector. \\[4pt]
\textbf{Equal-RMS carrier control.} At second order, equal local RMS
polarization gives the same steady source when the carriers see the same DC
coefficient. TI also supplies a beat; either waveform can induce an oscillation
if the mean shift crosses a bifurcation.
& Demonstrated in the PING mean field (Fig.~\ref{fig:qif_regimes}c,d).
The corresponding spiking-network and in vivo controls are not provided here.
& Compare two tones with one carrier of matched RMS polarization and duration,
checking the DC coefficient across the carrier interval. Measure mean activity,
intrinsic oscillations, and beat lock-in separately to distinguish the steady
shift from the beat response. \\
\end{longtable}
\endgroup
